% ==== P2 draft O1 — §2.2 대체 블록 (L106–L219 대체) ====
\subsection{단일 삽입 자리 --- 점유 통계와 그 요동}\label{sec:sm-site}
\textbf{(a) 출발 --- 두 가정과 원형 미시상태.} 격자의 삽입 자리 \emph{하나}를 떼어 살펴본다.
자리는 전해질 너머의 Li 저장조와 입자를 교환하므로 식~\eqref{eq:sm-gc} 의 대정준 설정 그대로다. 유도에
앞서 두 가정을 명시한다.
\begin{itemize}[leftmargin=2.2em,itemsep=0.1em]
\item \textbf{가정 1 (hard-core 배타).} 한 자리에는 Li 이온이 $0$ 개 또는 $1$ 개만 들어간다. 근거는 양자
통계가 아니라 강한 단거리 반발이다 --- 자리의 부피와 이온 사이 정전$\cdot$입체 반발이 이중 점유의
에너지를 사실상 무한대로 올린다. 이 가정이 점유수 합을 $n\in\{0,1\}$ 로 자른다.
\item \textbf{가정 2 (자리 독립).} 이 소절에서는 자리의 에너지가 이웃 자리의 점유와 무관하다고 둔다.
이 가정은 \S\ref{sec:sm-mf} 에서 상호작용 파라미터 $\Omega$ 로 완화된다.
\end{itemize}
표준 이론이 정확히 이 자리에 있다. 먼저 점유 상태를 우물 바닥 에너지 $\varepsilon_0$ 하나로 대표하는
원형을 세우고, 이어 점유 이온의 진동 자유도를 켜서 넓힌다. 원형의 미시상태는 둘이다 --- 빈 자리($n=0$,
에너지 $0$ 기준)와 점유 자리($n=1$, 에너지 $\varepsilon_0$).

\textbf{(b) 연산 --- 두 상태의 대정준 합.} 식~\eqref{eq:sm-gc} 의 대정준 합에 이 두 미시상태를 넣으면, 각
상태에 Gibbs 인자 $e^{-\beta(\varepsilon_0-\mu)n}$($\beta=1/k_BT$)을 곱해 더한 것이다:
\[
\Xi_1^{0}\;=\;\sum_{n=0,1}e^{-\beta(\varepsilon_0-\mu)n}
\;=\;1+e^{-\beta(\varepsilon_0-\mu)}.
\]
합의 결과 $\Xi_1^{0}$ 은 자리 1개의 \emph{대정준} 분배함수로, \S\ref{sec:sm-ensemble} 의 canonical $Z$ 와는
앙상블이 다르므로 기호부터 구분해 적는다(첨자 $1$ $=$ 자리 1개, 위첨자 $0$ $=$ 내부 자유도를 끈 원형).

\textbf{(c) 중간식 --- 원형의 평균 점유.} 자리의 평균 점유는 두 상태의 확률 가중 평균이다. 분자에 점유
상태의 Gibbs 인자를, 분모에 $\Xi_1^{0}$ 을 놓아 정리하면
\[
\langle n\rangle^{0}
=\frac{0\cdot1+1\cdot e^{-\beta(\varepsilon_0-\mu)}}{\Xi_1^{0}}
=\frac{e^{-\beta(\varepsilon_0-\mu)}}{1+e^{-\beta(\varepsilon_0-\mu)}}
=\frac{1}{1+e^{+\beta(\varepsilon_0-\mu)}},
\]
마지막 등식은 분자$\cdot$분모를 $e^{-\beta(\varepsilon_0-\mu)}$ 로 나눈 것이다.

\textbf{(d) 박스 --- 국소화 흡착의 점유 함수.} 두 결과를 모으면
\begin{equation}
\boxed{\;\Xi_1^{0}=1+e^{-\beta(\varepsilon_0-\mu)}\;,\qquad
\langle n\rangle^{0}=\frac{1}{1+e^{+\beta(\varepsilon_0-\mu)}}\;.}
\label{eq:sm-bare}
\end{equation}
이 두-상태 점유 함수가 국소화 흡착의 Langmuir 등온선이고 \cite{hill1960,fowler1939}, 삽입 전극에 적용한
원형은 McKinnon--Haering \cite{mckinnon1983} 이다. 점유 함수 \eqref{eq:sm-bare} 는 구조상 Fermi--Dirac
함수 $1/(e^{\beta(E-\mu)}+1)$ 와 이미 동형이다 --- 단일 준위 에너지가 $\varepsilon_0-\mu$ 인 셈이며,
공유하는 것은 ``한 상태에 입자 $0$ 또는 $1$''이라는 셈법 구조다. 그 셈법이 왜 페르미온 통계와 같은 꼴을
낳는지, 그리고 물리적으로 무엇이 다른지는 다음 배경 상자가 짚는다.

\begin{bgbox}
\textbf{페르미온·보손과 양자 통계 ---} 양자역학에서 같은 종류의 입자는 원리적으로 구별되지 않는다 ---
입자 하나하나에 꼬리표를 붙일 방법이 없다. 이 구별불가능성이 여러 입자의 파동함수를 입자 교환에 대해 두
부류로 가른다: 교환에 대칭인 것과 반대칭인 것. 대칭 파동함수를 갖는 입자가 \emph{보손}(boson)이며, 정수
스핀을 가지고 한 단일입자 상태에 몇 개든 들어갈 수 있어 Bose--Einstein 분포를 따른다. 반대칭 파동함수를
갖는 입자가 \emph{페르미온}(fermion)이며, 반정수 스핀을 가지고 Pauli 배타 원리로 한 준위의 점유수가
$0$ 또는 $1$ 로 제한되어 Fermi--Dirac 분포 $f(E)=1/(e^{\beta(E-\mu)}+1)$ 를 따른다. 페르미온의 $0/1$ 은
파동함수의 반대칭성이 강제하는 양자적 배타다. 반면 격자의 삽입 자리에서 점유수가 $0$ 또는 $1$ 인 것은 양자
배타가 아니라 기하 --- 자리의 부피와 이온 사이 정전$\cdot$입체 반발이 이중 점유를 금지한다. 두 경우의 물리는
다르지만 ``한 상태에 $0$ 또는 $1$''이라는 셈법의 구조가 같으므로, 같은 logistic(Fermi--Dirac) 대수형이
나온다. 진짜 양자 통계는 Chapter 2 에서 쓰인다 --- 격자 진동의 양자인 포논은 보손이라 Bose--Einstein
분포로, 전도 전자는 페르미온이라 Fermi--Dirac 분포로 각각의 엔트로피를 준다
\cite{mcquarrie1976,ashcroftmermin1976}.
\end{bgbox}

\textbf{확장 --- 점유 이온의 내부 자유도.} 점유된 이온은 우물 안 한 점에 얼어붙지 않는다 --- $n=1$ 인
자리에서 이온은 우물 안에서 진동하는 연속 자유도(위치$\cdot$운동량)를 가진다. 따라서 미시상태의 온전한
목록은 ``빈 자리($n=0$, 에너지 $0$ 기준)'' 하나와, ``점유 자리($n=1$, 우물 바닥 에너지 $\varepsilon_0$)에
내부 배치 $\Gamma$ 가 붙은 연속족''이다. 원형의 두 상태 합에서 점유 항의 무게를 이 내부 배치에 대한
적분으로 바꾸면
\begin{equation}
\Xi_1\;=\;\sum_{n=0,1}e^{\beta\mu n}\,z_n
\;=\;1+q(T)\,e^{-\beta(\varepsilon_0-\mu)},
\qquad
q(T)\equiv\int e^{-\beta H_\mathrm{int}(\Gamma)}\,\frac{\dd\Gamma}{h^3}
\label{eq:partfn}
\end{equation}
--- $z_n$ 은 점유수 $n$ 인 상태들의 canonical 분배함수로, $z_0=1$(빈 자리는 상태 하나),
$z_1=e^{-\beta\varepsilon_0}q(T)$(우물 바닥 인자에 내부 자유도의 적분 $q(T)$ 가 곱해진다)이다. 내부
자유도를 켜면 점유 상태의 canonical 무게가 $1$ 에서 $q(T)$ 로 바뀌어, 위첨자를 뗀 완전한 단일 자리
대정준 분배함수 $\Xi_1$ 가 된다. Chapter 2 의 단일 자리 분배함수(그 문건 식~(eq:Z1))도 같은 기호 $\Xi_1$
로 통일하며, 그 문건은 우물 바닥 에너지를 처음부터 유효값 $\tilde\varepsilon$ 로 읽어 $q(T)$ 흡수를 마친
형태로 적는다. $q(T)$ 는 점유된 이온의 진동 등 국소 자유도의 단일 자리 분배함수다\footnote{이
$q(T)$(단일 자리 내부 분배함수)는 N6--N9 의 정규화 용량 좌표 $q{=}Q/Q_\cell$ 와 다른 양이다 --- 전자는
Part 0 의 국소 열역학량, 후자는 곡선 좌표.}. 일반적인 정의는 상태 합
$q(T)=\mathrm{Tr}\,e^{-\beta H_\mathrm{int}}$ 이고 식~\eqref{eq:partfn} 의 위상공간 적분은 그 고전
극한이다 --- 자리 우물을 3차원 조화 퍼텐셜(진동수 $\omega_0$)로 근사하면 고전 극한에서
$q=(k_BT/\hbar\omega_0)^3$ 이고, 등방 가정(세 방향이 같은 $\omega_0$)을 풀어 데카르트 축마다 다른
진동수를 허용하는 일반형으로 넓히면 양자 상태 합으로는
$q=\prod_{i=1}^{3}[2\sinh(\beta\hbar\omega_i/2)]^{-1}$(축 지표 $i=1,2,3$ 마다 진동수 $\omega_i$)로 닫힌다.

\textbf{유효 자리 에너지 $\tilde\varepsilon$ 로의 흡수.} 점유 상태들의 확률 합이 곧 평균 점유다:
\begin{equation}
\langle n\rangle=0\cdot P_0+1\cdot P_1
=\frac{q(T)\,e^{-\beta(\varepsilon_0-\mu)}}{1+q(T)\,e^{-\beta(\varepsilon_0-\mu)}},
\label{eq:sm-occmid}
\end{equation}
그리고 식~\eqref{eq:sm-gc} 의 셋째 식으로 교차검증하면 $k_BT\,\partial\ln\Xi_1/\partial\mu$ 가 같은
결과를 준다. 여기서 적분이 남긴 계수 $q(T)$ 는 \emph{분자와 분모의 점유-상태 항 양쪽에 같은 인수로}
붙어 있다. 따라서 결과는 $q$ 와 $\varepsilon_0$ 을 따로 아는 데 의존하지 않고 한 조합에만 의존하며,
그 조합은 지수 안으로 흡수된다:
\begin{equation}
q(T)\,e^{-\beta\varepsilon_0}\;=\;e^{-\beta\tilde\varepsilon(T)},
\qquad
\tilde\varepsilon(T)\;\equiv\;\varepsilon_0-k_BT\ln q(T)
\label{eq:sm-epstilde}
\end{equation}
--- $\tilde\varepsilon$ 는 점유 자리의 \emph{자유에너지}(우물 바닥 에너지 $+$ 내부 자유도의 자유에너지
$-k_BT\ln q$)다. 자리 점유의 자유에너지 비용을 $\Delta\mu\equiv\tilde\varepsilon-\mu_\mathrm{Li}$ 로
줄여 쓰면 식~\eqref{eq:sm-occmid} 는 $e^{-\beta\Delta\mu}/(1+e^{-\beta\Delta\mu})$ 가 되어, 내부
자유도가 있어도 함수형은 두-상태 점유의 logistic 그대로 보존된다.

\textbf{박스 --- 점유율 $\theta$.} 분모$\cdot$분자를 $e^{-\beta\Delta\mu}$ 로 나누면
\begin{equation}
\boxed{\;\theta\;\equiv\;\langle n\rangle\;=\;\frac{1}{1+e^{+\beta\,\Delta\mu}}\;,
\qquad \Delta\mu\equiv\tilde\varepsilon(T)-\mu_\mathrm{Li}\;.}
\label{eq:fermifn}
\end{equation}
\begin{verifybox}
극한 검산 --- $\beta\to\infty$($T\to0$)에서 $\Delta\mu\gtrless0$ 에 따라 $\theta\to0/1$ 의 계단
($T\to0$ 에서 $\tilde\varepsilon$ 는 점유 상태의 바닥 에너지로 수렴한다 --- 고전 $q$ 면 $k_BT\ln q\to0$
이라 $\varepsilon_0$, 양자 조화 $q$ 면 영점 에너지를 더한 $\varepsilon_0+\tfrac32\hbar\omega_0$ --- 어느
쪽이든 상수라 계단 구조는 동일),
$\beta\to0$($T\to\infty$)에서는 $\beta\Delta\mu\to-\ln q$ 이므로 내부 자유도가 없는 기준($q\equiv1$)에서
$\theta\to\tfrac12$(두 상태 등확률)이고, $q(T)$ 가 온도와 함께 증가하는 실제 조화 우물에서는 점유 상태의
위상공간 증가가 이겨 $\theta\to1$ 로 향한다; $\mu_\mathrm{Li}\to\pm\infty$ 에서
$\theta\to1/0$; 내부 자유도가 없는 극한 $q\to1$ 에서 $\tilde\varepsilon\to\varepsilon_0$ 으로
식~\eqref{eq:sm-bare} 의 원형 박스를 그대로 회수한다. ★함수형 가드 --- 식~\eqref{eq:fermifn} 은
Fermi--Dirac 함수와 \emph{동형}이지만 물리량은 다르다: 공유하는 것은 ``한 자리에 입자 $0$ 또는 $1$''이라는
배타 점유의 대수 구조뿐이고(가정 1), 여기의 입자는 전자가 아니라 Li 이온이며 배타성의 근거도 양자 통계가
아니라 자리 하나에 이온 하나라는 기하다 --- 페르미온의 $0/1$ 이 파동함수 반대칭성(Pauli 배타)에서 오는
것과 달리, 여기 $0/1$ 은 자리 기하에서 온다.
\end{verifybox}

한 걸음 더 --- $n\in\{0,1\}$ 라 $n^2=n$ 이므로 점유의 요동(분산)이 닫힌 꼴로 나온다:
\begin{equation}
\mathrm{var}(n)=\langle n^2\rangle-\langle n\rangle^2=\theta(1-\theta),
\qquad
\frac{\partial\theta}{\partial(\beta\mu)}=\theta(1-\theta)
\label{eq:sm-flucres}
\end{equation}
(둘째 식은 식~\eqref{eq:fermifn} 의 $\beta\mu$ 미분). \emph{점유의 $\mu$-감도 $=$ 점유 요동} --- 요동--응답
(fluctuation--response) 관계의 최소 사례이고, 본론 평형 peak 의 종 $\xi_\eq(1-\xi_\eq)$(\S\ref{sec:eqpeak})의
통계적 정체가 바로 이 분산이다.

유효 자리 자유에너지 $\tilde\varepsilon(T)$ 가 하류에서 가는 곳은 두 갈래다. 첫째, 온도를 고정하고 보면
$\tilde\varepsilon$ 는 상수 하나이고, \S\ref{sec:sm-electro} 의 전기화학 결선에서 전이 중심 $U_j$ 로
흡수된다 --- 그 소절과 본론이 쓰는 몰당 기준 에너지 $\mu^0$ 는 정확히 $\tilde\varepsilon$ 의 몰 환산
$N_A\tilde\varepsilon$ 이다. 둘째, 온도를 미분하면 내부 자유도의 자유에너지 몫
$f_\mathrm{int}=-k_BT\ln q$ 가 자리당 엔트로피
\begin{equation}
s_\mathrm{int}\;=\;-\frac{\partial f_\mathrm{int}}{\partial T}
\;=\;k_B\,\frac{\partial\,(T\ln q)}{\partial T}
\label{eq:sm-sint}
\end{equation}
를 주는데, 이것이 Chapter 2 의 vibrational 엔트로피 사슬의 단일 자리 원형이다 --- 여기 자리 하나의
세 축 진동수 $\omega_i$($i=1,2,3$)를 격자 전체의 모든 진동 모드를 헤아리는 지표 $k$ 로 넓혀, 조화 모드의
$q_k=[1-e^{-\beta\hbar\omega_k}]^{-1}$ 를 식~\eqref{eq:sm-sint} 에 넣으면 그 문건의 모드당 엔트로피
$s_k=k_B[(1+\bar n_k)\ln(1+\bar n_k)-\bar n_k\ln\bar n_k]$(Bose--Einstein 들뜸수 $\bar n_k$)가 그대로
나온다. 기준점에 관한 사소한 주의로, 확장에서 든 $[2\sinh(\beta\hbar\omega/2)]^{-1}$ 은 우물 바닥 기준
(영점 에너지를 $q$ 에 포함)이고 여기의 $[1-e^{-\beta\hbar\omega}]^{-1}$ 은 영점 에너지를
$\varepsilon_0$ 쪽에 포함시킨 기준이다 --- 두 규약의 $q$ 는 인자 $e^{-\beta\hbar\omega/2}$(영점 에너지
$\hbar\omega/2$ 의 Boltzmann 인자)만큼 다른데, 이 인자가 $T\ln q$ 에 남기는 기여는 온도 무관 상수
$-\hbar\omega/(2k_B)$ 이므로 온도 미분, 곧 엔트로피는 동일하다. 따라서 전이 중심의 온도 기울기
$\dd U_j/\dd T=\Delta S_{\rxn,j}/F$(\S\ref{sec:center})에서 진동
엔트로피가 차지하는 몫의 미시적 기원이 바로 $\tilde\varepsilon(T)$ 의 온도 의존, 곧 $q(T)$ 다.
% ==== 블록 끝 ====
% NOTE:
% Read Coverage:
%   - _sections/ch1_sec02a_part0.tex 전문 L1–269 (대상 원문, 특히 L106–219)
%   - results/V1020_STYLE_RUBRIC.md 전문 L1–51
%   - _sections/ch1_sec02b_part0.tex L1–60 (§2.4 가 받는 μ(θ)·Ω 슬롯)
%   - _sections/ch2_sec01_partition.tex L16–50 (eq:Z1·eq:occ 원형, ε̃ 흡수·s_int 기대)
%   - _sections/ch1_preamble.tex L20–49 (bgbox=\newtheorem*{bgbox}{배경}·verifybox·매크로 확인)
% 보존 체크:
%   P-1 가정 1/2 명시: (a) 출발 itemize (가정 1 hard-core 배타·가정 2 자리 독립).
%   P-2 미시상태 온전 기술(점유 상태 연속 내부 배치): 확장 첫 단락 "온전한 목록 ... 내부 배치 Γ 가 붙은 연속족".
%   P-3 라벨 6종 최종 수식 불변: eq:partfn·eq:sm-occmid·eq:sm-epstilde·eq:fermifn·eq:sm-flucres·eq:sm-sint 모두 원문 수식 그대로. 신규 eq:sm-bare 추가.
%   P-4 q(T)↔정규화 용량 좌표 q 구분 각주: eq:partfn 뒤 q(T) 정의 문장의 \footnote (원문 그대로).
%   P-5 조화 우물 q 고전 (k_BT/ħω₀)³·양자 ∏[2sinh(βħω_i/2)]^{-1}: 확장 첫 단락 말미.
%   P-6 ★FD 동형 가드: verifybox 후반 "공유하는 것은 ~뿐이고 ... 여기의 입자는 ~ Li 이온 ... 기하다" (자족·bgbox 접속).
%   P-7 요동-응답 var(n)=θ(1-θ)=∂θ/∂(βμ): eq:sm-flucres + "μ-감도=요동" 단락.
%   P-8 ε̃ 하류 두 갈래(U_j 흡수·s_int=k_B∂(T ln q)/∂T·Ch2 vib 원형·영점 규약 주의): 마지막 두 단락, eq:sm-sint 포함.
%   P-9 grand canonical 문맥(eq:sm-gc 셋째 식 교차검증): 확장 "유효 자리 에너지" 단락 "k_BT ∂ln Ξ₁/∂μ 가 같은 결과" + (b) eq:sm-gc 대정준 합.
%   P-10 Ch2 Ξ₁ 통일 문장: eq:partfn 뒤 "Chapter 2 의 단일 자리 분배함수(그 문건 식 (eq:Z1))도 같은 기호 Ξ₁ 로 통일".
% 신규 라벨: eq:sm-bare (원형 두-상태 박스: Ξ₁⁰·⟨n⟩⁰). 기타 라벨 전부 기존 것 재사용(신규 없음).
% 자체검수:
%   - 극한: β→∞ 계단, β→0 시 q≡1→θ→½·실제 q(T)→θ→1, μ→±∞→θ→1/0, q→1 시 eq:sm-bare 원형 박스 회수 명시. Ξ₁⁰=1+e^{-β(ε₀-μ)} 는 Ch2 eq:Z1 과 형태 일치(정합).
%   - 부호: Gibbs 인자 e^{-β(ε₀-μ)n} → ⟨n⟩⁰=1/(1+e^{+β(ε₀-μ)}) 부호 일치; ε̃=ε₀-k_BT ln q 로 logistic 함수형 불변.
%   - rubric: D7(원형 박스→확장 동기→확장 박스→원형 회수 검산) 준수, 원형·확장 박스 분리; A2/D1 자기개정 서술 0(“원형→확장”은 물리적 층위 서술); FD 가드는 A2 예외 대조형; C3 페르미온(fermion)·보손(boson) 병기; 인용 화이트리스트만(hill1960·fowler1939·mckinnon1983·mcquarrie1976·ashcroftmermin1976).
